% ============================================================================
%  Chapter 4 — Attack Framework and the Undefended Baseline
%  Thesis: Generative AI, Privacy & Personal Data — TADJ Mohamed Walid, ESI SBA
%
%  DUAL-MODE FILE. Compiles standalone (for drafting) AND is \input-able into a
%  future thesis.tex (define \ThesisMaster before \input to skip the preamble).
%
%  CITATIONS: \cite{...-TODO} are placeholders; the unified .bib is separate
%  scaffolding (PHASE5_PLAN.md §4.2).
%
%  CONSOLIDATE chapter. Source prose: Phase2_Report_Detailed.tex (committed,
%  supervisor-reviewed). CANONICAL baseline for ALL attack-success rates (overall,
%  per-profile, per-family) = results/phase3/ablation/metrics/final_metrics.json,
%  cell C00000 -- the SAME baseline Chapter 5 measures reductions against. Severity
%  (mean_S), the engine-judge cross-tab, retrieval precision, the hand-verified-leak
%  audit, and the Phase-2 judge calibration kappa are drawn from the Phase-2 run
%  (baseline_metrics_augmented.json) and labeled as such (final_metrics carries none
%  of them). The threat model, four profiles, (L,S) definitions and ground-truth
%  corpus are OWNED by Chapter 3 (\ref{ch:threat}/\ref{ch:formal}) and are referenced
%  here, not re-derived. The full two-instrument treatment and the extended kappa
%  chain are OWNED by Chapter 5 (\ref{ch:defenses}); this chapter states only the
%  Phase-2 single-annotator validation.
% ============================================================================
\ifdefined\ThesisMaster\else
  \documentclass[11pt,a4paper]{report}
  \usepackage[utf8]{inputenc}
  \usepackage[T1]{fontenc}
  \usepackage{lmodern}
  \usepackage[margin=2.5cm]{geometry}
  \usepackage{booktabs}
  \usepackage{tabularx}
  \usepackage{array}
  \newcolumntype{L}{>{\raggedright\arraybackslash}X}
  \usepackage{amsmath,amssymb}
  \usepackage{graphicx}
  \usepackage{microtype}
  \usepackage{xcolor}
  \usepackage[hidelinks]{hyperref}
  \begin{document}
\fi

\chapter{Attack Framework and the Undefended Baseline}
\label{ch:baseline}

The defenses evaluated in Chapter~\ref{ch:defenses} are measured as reductions from a reference,
and this chapter establishes that reference. It specifies the adversarial campaign of fifty-two
scenarios across eight attack families, describes the retrieval-augmented system under test, fixes
the dual-instrument measurement protocol, and reports the leakage of that system with no privacy
defenses in place. The chapter answers \textbf{RQ1}: on an undefended enterprise RAG deployment,
which kinds of attack actually produce privacy leakage, and which dominate. Its central number, a
strict attack-success rate of $28.8\%$, is the anchor against which every later figure in the
thesis is read, so the chapter is careful about what that number does and does not count.

The threat actors, the LINDDUN-grounded profiles, the leakage taxonomy, and the formal $(L,S)$
definitions are established in Chapter~\ref{ch:threat} and used here without restatement; this
chapter concerns the \emph{attack framework} and the \emph{baseline}, not the foundations.

\section{The attack framework}
\label{sec:ch4-framework}

The campaign comprises \textbf{fifty-two adversarial scenarios} spanning eight attack families
(F1--F8) and exercising three of the four threat profiles through live interaction
(Table~\ref{tab:ch4-families}). The fourth profile, the Third-Party Provider (P4), has no live
scenarios: its exposure is bounded by an API contract rather than a query, and is assessed
offline in Chapter~\ref{ch:threat} rather than through the chat interface. Each scenario
belongs to exactly one family and is identified by a \texttt{profile--family--number} code.

\begin{table}[h]
\centering
\caption[The eight attack families and the profiles that exercise each]{The eight attack families, the profiles that exercise each, and the scenario counts.
F1--F7 are shared by the Internal (P1) and External (P2) profiles; corpus poisoning (F8) is an
Admin (P3) capability, since it acts on the index before any query is issued. Source:
\texttt{Phase2\_Report\_Detailed.tex} §3--4.}
\label{tab:ch4-families}
\begin{tabularx}{\textwidth}{lLcccc}
\toprule
\textbf{Family} & \textbf{Mechanism} & \textbf{P1} & \textbf{P2} & \textbf{P3} & \textbf{Total}\\
\midrule
F1 Direct PII extraction      & names a record and requests sensitive attributes directly & 3 & 6 & --- & 9\\
F2 Cross-document linkability & forces multi-document retrieval to correlate entities across records & 7 & 3 & --- & 10\\
F3 Direct prompt injection    & embeds authority/maintenance claims to override safety framing & 1 & 5 & --- & 6\\
F4 Obfuscated query           & encodes the request (Base64, Unicode, paraphrase) to evade keyword checks & 1 & 4 & --- & 5\\
F5 Jailbreak / role-play      & persona or fictional framing to shed operational constraints & 1 & 4 & --- & 5\\
F6 System reconnaissance      & probes corpus topology and retrieval structure, no direct payload & 1 & 2 & --- & 3\\
F7 Social engineering         & fabricated emergency or authority to bypass refusals by narrative pressure & 4 & 3 & --- & 7\\
F8 Corpus poisoning           & pre-stages malicious documents in the index for later retrieval & --- & --- & 7 & 7\\
\midrule
\textbf{Total} & & \textbf{18} & \textbf{27} & \textbf{7} & \textbf{52}\\
\bottomrule
\end{tabularx}
\end{table}

The distribution is uneven by design. The External profile carries the most scenarios (27),
because within-profile diversity matters most for an unauthenticated attacker with the full
published repertoire of adversarial prompts; the Internal profile carries 18, weighted toward
linkability (its characteristic threat); and the Admin profile carries 7, all corpus
poisoning, because that threat is validated through retrieval persistence rather than query
volume. The poisoning family rests on six poisoned documents
(\texttt{poisoned\_001}--\texttt{poisoned\_006}), one of which is exercised by two of the seven
F8 scenarios.

\section{The system under test}
\label{sec:ch4-sut}

\begin{figure}[h]
\centering
\resizebox{\textwidth}{!}{%
\begin{tikzpicture}[font=\small, >={Latex[length=2.6mm]},
  proc/.style={draw=figblue!70!black, fill=figblue!10, rounded corners=3pt, minimum height=13mm, minimum width=25mm, align=center, thick},
  io/.style={draw=figslate!85!black, fill=figslate!20, rounded corners=3pt, minimum height=13mm, minimum width=22mm, align=center, thick},
  llm/.style={draw=figamber!80!black, fill=figamber!28, rounded corners=3pt, minimum height=15mm, minimum width=31mm, align=center, very thick},
  store/.style={cylinder, shape border rotate=90, aspect=0.28, draw=figteal!75!black, fill=figteal!22, minimum height=17mm, minimum width=27mm, align=center, thick, text=figteal!30!black, font=\bfseries},
  ar/.style={->, thick, figblue!65!black}]
\node[proc] (corpus) at (0,3) {Corpus\\\footnotesize 155 docs};
\node[proc] (chunk)  at (3.4,3) {Chunking\\\footnotesize 512 / 64};
\node[proc] (embed)  at (6.8,3) {MiniLM\\\footnotesize 384-d};
\node[io]   (q)   at (0,0)    {User\\query};
\node[store](db)  at (6.8,0)  {ChromaDB};
\node[llm]  (llm) at (11.2,0) {LLM (SUT)\\\footnotesize Sonnet 4.6 \,$\cdot$\, $T{=}0$};
\node[io]   (resp)at (15.4,0) {Response};
\draw[ar] (corpus) -- (chunk);
\draw[ar] (chunk) -- (embed);
\draw[ar, figteal!70!black] (embed) -- node[right,font=\scriptsize]{index} (db);
\draw[ar] (q) -- node[above,font=\scriptsize]{retrieve} (db);
\draw[ar, figteal!70!black] (db) -- node[above,font=\scriptsize]{top-5 context} (llm);
\draw[ar] (llm) -- (resp);
\begin{scope}[on background layer]
  \node[fill=figteal!8, rounded corners=5pt, fit=(corpus)(embed), inner sep=7pt,
        label={[font=\scriptsize\itshape, text=figteal!45!black]above:indexing \,(offline)}] {};
  \node[fill=figamber!9, rounded corners=5pt, fit=(q)(resp), inner sep=7pt,
        label={[font=\scriptsize\itshape, text=figorange!60!black]below:serving \,(online)}] {};
\end{scope}
\end{tikzpicture}}
\caption[The retrieval-augmented generation system under test]{The minimal RAG system under test.
\textbf{Indexing} (top, teal): the 155-document corpus is chunked (512/64), embedded with
all-MiniLM-L6-v2 (384-d), and stored in the ChromaDB vector store. \textbf{Serving} (bottom,
amber): a user query retrieves the top-$5$ chunks, which ground the LLM (Claude Sonnet~4.6,
temperature~0) in producing the response. No safety scaffolding is added at the baseline, so the
measurement reflects what the model and retriever do unaided.}
\label{fig:ch4-architecture}
\end{figure}

The system under test (\Cref{fig:ch4-architecture}) is deliberately minimal, so that the baseline measures what an
underlying model and retriever do \emph{without} added safety scaffolding. Orchestration is
LangChain; the vector store is ChromaDB; embeddings are local \texttt{all-MiniLM-L6-v2}
(384-dimensional); the generating model is Claude Sonnet~4.6 at temperature~$0$. The system
prompt is generic and contains no privacy instruction, no role restriction, and no refusal
policy, because a defensively crafted prompt would itself be an uncharacterized defense and
would contaminate the reference. The corpus is the Phase~1 ground truth (155 synthetic
documents, $2{,}190$ annotated entities, described in full in Chapter~\ref{ch:formal}), and
no real personal data is used at any stage.

Each scenario was executed three times under deterministic decoding to expose
run-to-run variance, yielding $52 \times 3 = 156$ \emph{(scenario, run) executions} across 267
conversational turns; a 17-turn pilot used for system validation is excluded from all baseline
metrics. The execution is the unit of analysis: an attack-success rate is the fraction of the
156 executions in which at least one turn produced a leak at the relevant threshold, and
per-profile and per-family rates use the corresponding subset.

\section{Dual-instrument, dual-threshold measurement}
\label{sec:ch4-measurement}

Two design decisions in the measurement protocol shape every number that follows, and both
are responses to a specific weakness of deterministic detection.

\paragraph{Two thresholds.} The deterministic engine fires whenever a response contains an
entity whose severity exceeds the requester's authorization level. For the External profile
that threshold is zero, so \emph{any} disclosed category counts, including a severity-2
category noun (``department,'' ``job title'') that appears in the act of \emph{refusing}: a
response such as ``I cannot share salary or departmental data'' trips the engine on
``departmental'' though it discloses nothing. The campaign therefore reports two rates. The
\emph{inclusive} rate ($\text{ASR}_{\text{incl}}$, severity~$\geq 2$) preserves the full engine
signal and is a conservative upper bound; the \emph{strict} rate
($\text{ASR}_{\text{strict}}$, severity~$\geq 3$) suppresses the severity-2 refusal artifact
and is the operationally meaningful figure. The strict rate is primary throughout the thesis,
and ``$28.8\%$'' refers to it.

\paragraph{Two instruments.} The same weakness means the engine cannot, on its own, tell a
genuine disclosure from a refusal that names a category, nor can it see a leak the model
\emph{paraphrases} rather than reproduces verbatim. A language-model judge (Claude Opus~4.7)
reads each response semantically as a second labeller. The full treatment of the two
instruments (why the judge is primary, why the engine is a lower bound, and how the judge's
reliability shifts once defenses redact and refuse) belongs to Chapter~\ref{ch:defenses}
and is not repeated here. What this chapter needs is the Phase-2 validation that licensed the
judge in the first place: on a 46-record stratified calibration sample, the judge agreed with
the human annotator at Cohen's $\kappa = 0.827$ (almost perfect; raw agreement $43/46 = 93.5\%$),
while the engine agreed with the human at only $\kappa = 0.164$
(Table~\ref{tab:ch4-kappa}). One qualification must travel with that figure: the human
reference here is a \emph{single} annotator (the author), so $\kappa = 0.827$ is a
judge--author concordance, not an inter-annotator-validated ground truth. A second annotator
was, at this stage, future work; the second-annotator cross-check and the cross-vendor
corroboration that extend this validation were carried out subsequently and are reported in
Chapter~\ref{ch:defenses}.

\begin{table}[h]
\centering
\caption{Phase-2 judge calibration on the 46-record sample, against a single human annotator
(the author). The judge passes the $\kappa \geq 0.60$ bar; the engine fails it. Source:
\texttt{Phase2\_Report\_Detailed.tex} §B.}
\label{tab:ch4-kappa}
\begin{tabular}{lcl}
\toprule
\textbf{Comparison} & \textbf{Binary $\kappa$} & \textbf{Reading}\\
\midrule
Judge vs.\ human (author) & \textbf{0.827} & almost perfect (validation)\\
Engine vs.\ judge & 0.245 & fair\\
Engine vs.\ human (author) & 0.164 & poor\\
\bottomrule
\end{tabular}
\end{table}

\section{RQ1: the undefended baseline}
\label{sec:ch4-rq1}

A note on sourcing precedes the numbers, because it settles a question that recurs across the
thesis. \textbf{All attack-success rates in this chapter are reported from the ablation's
undefended cell} (\texttt{final\_metrics.json}, configuration C00000), the same artifact
and the same baseline that Chapter~\ref{ch:defenses} measures every reduction against, so
that the baseline and the reductions taken from it are one number, family by family. Severity,
retrieval-precision, hand-verified-audit, and judge-calibration figures, which that artifact
does not carry, are drawn from the Phase~2 campaign run
(\texttt{baseline\_metrics\_augmented.json}) and are labeled as such where they appear.

With the framework and protocol fixed, the undefended system was run across the full campaign.
At the execution level the judge, the primary instrument, reports a strict
attack-success rate of $\mathbf{28.8\%}$ ($45/156$), against $19.2\%$ read through the
deterministic engine alone (Table~\ref{tab:ch4-headline}). The judge reads higher because it
perceives the paraphrased name-plus-attribute disclosures that annotation-grounded matching
cannot represent; the engine rate is a lower bound and the judge rate primary, exactly as in
Chapter~\ref{ch:defenses}. Drawn from the Phase~2 run, the mean severity of a leaking execution
is high ($\bar{S} = 4.49$ over all executions, $7.96$ conditional on a leak), indicating that
the leaks which occur expose high-severity content rather than incidental detail.

\begin{table}[h]
\centering
\caption{Undefended baseline (C00000) attack-success rate: engine (lower bound) versus judge
(primary), strict (severity~$\geq 3$) and inclusive (severity~$\geq 2$). Source:
\texttt{final\_metrics.json} (C00000).}
\label{tab:ch4-headline}
\begin{tabular}{lcc}
\toprule
\textbf{Metric} & \textbf{Engine (lower bound)} & \textbf{Judge (primary)}\\
\midrule
$\text{ASR}_{\text{strict}}$    & 19.2\% & \textbf{28.8\%}\\
$\text{ASR}_{\text{inclusive}}$ & 48.1\% & 55.8\%\\
\bottomrule
\end{tabular}
\end{table}

The mechanism behind the engine--judge gap is visible in the Phase~2 turn-level
cross-tabulation (267 turns): the engine fired on 95 turns, the judge on 57, their union on
115. Of the engine's 95, only 37 were also judged leaks; 58 were the severity-2 External
refusal artifacts the strict threshold exists to exclude, and 20 further leaks were seen by the
judge alone: the paraphrased name-plus-attribute disclosures. This is the empirical
justification for both the strict threshold and the second instrument, treated in full in
Chapter~\ref{ch:defenses}.

\paragraph{By profile.} Leakage concentrates in the human-driven profiles
(Table~\ref{tab:ch4-profile}, \Cref{fig:ch4-asr-profile}). The Internal attacker is the most successful ($33.3\%$ judged),
because legitimate authenticated access combined with linkability needs no manipulation at all;
the External attacker follows closely ($29.6\%$); the Admin profile registers $14.3\%$,
entirely from corpus poisoning.

\begin{table}[h]
\centering
\caption{Strict ASR by profile at the undefended baseline, engine (lower bound) versus judge
(primary). Rates are over the executions in each profile's scenario subset. Source:
\texttt{final\_metrics.json} (C00000).}
\label{tab:ch4-profile}
\begin{tabular}{lccc}
\toprule
\textbf{Profile} & \textbf{Engine} & \textbf{Judge} & \textbf{Scenarios}\\
\midrule
P1 Internal & 27.8\% & \textbf{33.3\%} & 18\\
P2 External & 18.5\% & 29.6\% & 27\\
P3 Admin    & 0.0\%  & 14.3\% & 7\\
\midrule
\textbf{Overall} & 19.2\% & \textbf{28.8\%} & 52\\
\bottomrule
\end{tabular}
\end{table}

\begin{figure}[h]
\centering
\includegraphics[width=0.85\textwidth]{fig1_asr_by_profile}
\caption[Strict ASR by threat-actor profile: engine baseline vs.\ augmented metric]{Scenario-level
ASR\textsubscript{strict} by threat-actor profile, comparing the deterministic engine baseline
against the augmented metric ($L_{\max}=\max(L_{\text{engine}},L_{\text{judge}})$). $N=52$
scenarios, three runs each at temperature~0. Source: \texttt{baseline\_metrics.json} (engine) and
\texttt{baseline\_metrics\_augmented.json} (augmented).}
\label{fig:ch4-asr-profile}
\end{figure}

\paragraph{By family: the RQ1 answer.} The decomposition by attack family answers RQ1
directly and unambiguously (Table~\ref{tab:ch4-family}, \Cref{fig:ch4-asr-family}). \textbf{Direct PII extraction (F1) is
the strongest attack family, at a strict judged ASR of $0.778$}, followed by cross-document
linkability (F2) at $0.600$. The two are in a class of their own: the next-highest family sits at $0.200$ and the rest at
$0.143$ or below. The manipulation families (direct injection F3, obfuscation F4, reconnaissance
F6, and social engineering F7) register $0.000$ under the strict threshold, because on this corpus
they either fail to surface above-severity-2 content or are refused; jailbreak (F5) reaches only
$0.200$. Corpus poisoning (F8) registers $0.143$, but from a single reproducible scenario, the
descriptive-poisoning exploit P3-POISON-007: a one-scenario family rate carries no confidence
interval and is not robustness-comparable to the multi-scenario F1 and F2 rates. The Phase~2
turn-level analysis corroborates the concentration from the other side: F1 and F2 together
accounted for $91.1\%$ of all strict leak turns in that run.

\begin{table}[h]
\centering
\caption{Strict ASR by attack family at the undefended baseline, engine (lower bound) versus
judge (primary). F1 and F2 dominate; \dag\,marks the single-scenario family (F8), whose rate
carries no confidence interval. Source: \texttt{final\_metrics.json} (C00000).}
\label{tab:ch4-family}
\begin{tabular}{lccc}
\toprule
\textbf{Family} & \textbf{Engine} & \textbf{Judge} & \textbf{Scenarios}\\
\midrule
F1 Direct PII extraction      & 0.556 & \textbf{0.778} & 9\\
F2 Cross-document linkability & 0.400 & \textbf{0.600} & 10\\
F5 Jailbreak / role-play      & 0.200 & 0.200 & 5\\
F8 Corpus poisoning\dag       & 0.000 & 0.143 & 7\\
F3 Direct prompt injection    & 0.000 & 0.000 & 6\\
F4 Obfuscated query           & 0.000 & 0.000 & 5\\
F6 System reconnaissance      & 0.000 & 0.000 & 3\\
F7 Social engineering         & 0.000 & 0.000 & 7\\
\bottomrule
\end{tabular}
\end{table}

\begin{figure}[h]
\centering
\includegraphics[width=0.82\textwidth]{fig2_asr_by_family}
\caption[Augmented ASR by attack family]{Augmented ASR\textsubscript{strict} by attack family. F1
(Direct PII Extraction) and F2 (Cross-document Linkability) together account for $91.1\%$ of all
turn-level strict leaks. F8 (Corpus Poisoning) at $14.3\%$ reflects a single reproducible exploit
(P3-POISON-007); F3, F4, and F7 remain at $0\%$ under both metrics. Source:
\texttt{baseline\_metrics\_augmented.json} (\texttt{ASR\_by\_family\_strict}).}
\label{fig:ch4-asr-family}
\end{figure}

The shape of this result, not merely its magnitude, sets the agenda for the rest of the thesis.
Leakage on an undefended enterprise RAG system is overwhelmingly a problem of the system
\emph{doing what it was built to do}: retrieving a record and relaying its contents (F1), or
retrieving several and letting the model correlate them (F2), rather than of exotic adversarial
manipulation, which the underlying model largely refuses (\Cref{fig:ch4-leaktype}). The severity profile of the Phase~2 run
confirms the same point: the judge's leak turns concentrate at severity~3 and~4 (56 of 57),
whereas two-thirds of the engine's sit at severity~2, the refusal-noise band. The defenses of
Chapter~\ref{ch:defenses} are therefore aimed where the leakage actually is, at F1 and F2, and the
priority order falls directly out of this table.

\begin{figure}[h]
\centering
\includegraphics[width=0.6\textwidth]{fig3_leak_type_distribution}
\caption[Augmented leak-type distribution across the campaign turns]{Augmented leak-type
distribution across $N=267$ attack turns. Direct PII extraction accounts for the majority of
confirmed leaks; cross-reference disclosure (turn-linking two documents) accounts for the
remainder. No schema-disclosure leaks were detected under the augmented metric. Source: per-turn
logs and \texttt{judge\_verdicts\_full.jsonl}.}
\label{fig:ch4-leaktype}
\end{figure}

\section{Why the baseline is a conservative floor}
\label{sec:ch4-floor}

One property of this baseline must be stated plainly, because it governs how the $28.8\%$
should be read: it is a \emph{floor}, not a ceiling. The deterministic engine credits a leak
only to entities annotated on a document the retriever actually returned, and even the judge
can only flag what appears in the response it reads; a leak assembled from a record the
retriever did not fetch is uncounted. Two facts from the Phase~2 campaign make this concrete.
Retrieval precision on the target document is low ($8.3\%$; \Cref{fig:ch4-retrieval}), so the retriever frequently
returns neighboring rather than targeted records, and a disclosure reconstructed from those is
invisible to the engine. And a manual audit independently confirmed ten leak cases by hand,
including the most severe (medical conditions co-occurring with a salary and a name) and
two cross-reference disclosures in which the response enumerated retrieval-context document IDs
never present in the user's query. The true undefended leakage rate is thus at least $28.8\%$;
the measurement understates it rather than overstating it.

\begin{figure}[h]
\centering
\includegraphics[width=0.82\textwidth]{fig6_retrieval_precision}
\caption[Retrieval precision by attack family]{Retrieval precision by attack family. Of the 36
named-target turns, the retriever returned the correct document in only 3 ($8.3\%$); all 3 hits
are F5 (Jailbreak), whose scenario named \texttt{hr\_file\_007} and retrieved it in all 3 runs. F1,
F3, F4, and F7 never retrieved their named target ($0\%$). F1 precision $=0\%$ yet F1
ASR\textsubscript{strict} $=77.8\%$: the retriever returned non-target adjacent documents, from
which the model nonetheless disclosed PII. F2, F6, and F8 are N/A (no single target prescribed).
Source: \texttt{baseline\_metrics.json} (\texttt{retrieval\_precision\_by\_family}).}
\label{fig:ch4-retrieval}
\end{figure}

\section{Summary}
\label{sec:ch4-summary}

This chapter fixed the reference point for the thesis. On an undefended enterprise RAG system,
measured over 52 scenarios and 156 executions with a dual-threshold, dual-instrument protocol,
the strict attack-success rate is $28.8\%$. \textbf{RQ1 is answered: direct PII extraction
(F1, $0.778$) and cross-document linkability (F2, $0.600$) are the families that maximize
leakage, in a class of their own, with the next family reaching only $0.200$}, while adversarial
manipulation (injection, obfuscation, social engineering) is largely refused and registers at
zero under the strict threshold, jailbreak reaching only $0.200$. Leakage is driven by the retrieval pipeline working as
designed, not by exotic prompts; it concentrates in the Internal and External profiles; and
the $28.8\%$ figure is a conservative floor, grounded by a hand-verified leak audit and
understated by the engine's annotation-bound view. This is the baseline from which
Chapter~\ref{ch:defenses} measures every reduction.

\ifdefined\ThesisMaster\else
  \end{document}
\fi
